\chapter{Emerging Trends (2024--2026)}
\label{ch:emerging}

\epigraph{\itshape The cutting edge is where bibliographies become
incomplete.}{}

This chapter surveys directions that are still moving fast and likely
to shape the next phase of SciML. Citations indicate vintage 2024--2026.

\section{Foundation models for physics}
\paragraph{Multiple Physics Pretraining (MPP)~\cite{mccabe2024mpp}.}
Train a single transformer on diverse physical fields (compressible
NS, MHD, shear flows) using a unified channel-tokenised representation;
fine-tune to new tasks with little data.

\paragraph{Poseidon~\cite{herde2024poseidon}.}
Multiscale operator transformer pretrained on a heterogeneous PDE mix
with semi-group training; demonstrates strong transfer to unseen PDE
families and downstream sample efficiency.

\paragraph{DPOT~\cite{hao2024dpot}.}
Denoising pretrained operator transformer: combine masked autoencoding
with operator regression to scale across PDE corpora.

\paragraph{Other entrants.}
\begin{itemize}
\item \textbf{Aurora} (Microsoft, 2024): foundation atmospheric model for
weather, air quality, ocean.
\item \textbf{NeuralGCM~\cite{kochkov2024neuralgcm}}: differentiable
hybrid GCM with skill matching ECMWF on medium-range forecasts.
\item \textbf{GenCast~\cite{price2024gencast}}: diffusion ensemble
forecaster, top of WeatherBench 2.
\item \textbf{ORBIT (ORNL)}: 100-billion-parameter weather model on
Frontier.
\item \textbf{Pangu, GraphCast, FourCastNet}: 2022--2023 progenitors;
the field's first ``GPT moments''.
\end{itemize}

\paragraph{What is still missing.}
\begin{itemize}
\item Tokenisation of heterogeneous physical fields (different units,
symmetries, dimensionality).
\item Standardised evaluation across PDE families.
\item Truly equivariant foundation models at scale.
\item Long-horizon stability beyond a few hundred autoregressive steps.
\item Open data + open weights: still patchy outside weather/climate.
\end{itemize}

\section{Diffusion and flow-based generative simulators}
The 2025 wave fuses generative modelling with PDE solving:
\begin{itemize}
\item \textbf{PDE-Refiner, PDE-Diffusion}: iterative denoising as a
universal post-processor for operator outputs.
\item \textbf{SimDiffPDE~\cite{simdiffpde2025}}: plain transformer
diffusion as a competitive PDE solver across a wide problem set.
\item \textbf{PBFM~\cite{pbfm2025}}: physics-based flow matching
imposing PDE residuals into the flow-matching objective.
\item \textbf{GenCast}: probabilistic ensemble forecasting at
state-of-the-art skill.
\item \textbf{Diffusion priors with physics-guided
inference}~\cite{diffphys2026}: train task-agnostic priors, condition at
sampling time on PDE residuals --- decoupling priors from PDE
specifications.
\end{itemize}

\section{Transformer operators}
Transformers are the architecture backbone for the largest 2024--2026
SciML models:
\begin{itemize}
\item Attention naturally handles unstructured / variable-resolution
inputs.
\item Tokenisation of physical fields via patches, points, or learned
tokens.
\item Cross-attention between query coordinates and conditioning
fields (DeepONet-style branches).
\item Adaptive attention windows for multiscale problems.
\end{itemize}
Examples: Transolver, GAOT, OFormer, GNOT, MPP, DPOT, Poseidon.

\section{Rethinking PINNs}
A 2025 NeurIPS result~\cite{rethinking2025pinn} reframes PINN failure
modes as \emph{precision-induced optimisation pathologies} rather than
fundamental landscape problems --- L-BFGS' default convergence test
fires too early in FP32. Switching to FP64 and using a stricter
tolerance recovers vanilla-PINN convergence on benchmarks previously
considered unsolvable. This is reshaping how the community discusses
the limits of PINN training.

\section{Equation discovery at scale}
\begin{itemize}
\item \textbf{End-to-end transformer symbolic regression} (NeSymRes,
SymbolicGPT, OmegaPRM-style search) is closing on PySR for benchmark
problems.
\item \textbf{LLM-guided equation discovery}: large language models
proposing candidate symbolic forms, validated by numerical fitting
(LaSR, MathConstellation).
\item \textbf{Physics-aware program synthesis}: combining differentiable
physics with neural-symbolic search for governing equations, hybrid
ODEs, and stochastic models.
\end{itemize}

\section{AI for science programmes}
National-scale initiatives are reshaping funding and data:
\begin{itemize}
\item \textbf{DOE Foundation Models for Science} programmes (US):
billions of CPU/GPU hours, multi-physics datasets, public foundation
models.
\item \textbf{EU AI Factories / Destination Earth}: digital twins of
the planet with hybrid SciML.
\item \textbf{UK AISI / Royal Society reports} on AI in science.
\item \textbf{Industry labs}: DeepMind (GenCast, AlphaFold), NVIDIA
Modulus, Microsoft (Aurora, DiG, MatterGen), Meta FAIR (OC, ESM).
\end{itemize}

\section{Quantum-enhanced SciML}
Quantum-classical hybrid solvers (variational quantum eigensolvers,
quantum signal processing) intersect with neural surrogates for
near-term quantum chemistry and lattice QFT. Currently early-stage but
worth tracking.

\section{Active learning, autonomous experimentation, and AI agents}
LLM-powered agents that orchestrate simulation $+$ experiment $+$ ML
modelling are an emerging mode. Examples: A-Lab (materials), ChemCrow,
self-driving labs in materials science, agentic SciML pipelines that
plan, run, and analyse simulations end-to-end.

\section{Reinforcement learning for control of dynamical systems}
\begin{itemize}
\item \textbf{Plasma control} (TCV, ITER digital twins).
\item \textbf{Fluid control} (drag reduction, mixing).
\item \textbf{Climate engineering scenarios} (still controversial).
\item Differentiable physics + policy-gradient RL is converging into a
unified control framework.
\end{itemize}

\section{Where to bet next}
\begin{insight}{Five-year outlook}
\begin{itemize}
\item Foundation operators will become commodity; the moat moves to
\emph{data + downstream finetuning} pipelines.
\item Probabilistic / generative surrogates will overtake deterministic
ones for any high-stakes application requiring UQ.
\item Equivariant + symplectic + transformer hybrids are likely
architectural sweet spots.
\item Hybrid mechanistic + neural models (UDEs, NeuralGCM) will outpace
pure-NN approaches in safety-critical domains.
\item LLM-mediated equation discovery and agentic SciML pipelines will
democratise method design.
\end{itemize}
\end{insight}
